% F081: precise mathematical teaching diagram.
\begin{tikzpicture}[x=1cm,y=1cm,>=stealth,font=\small]
\draw[rounded corners=3pt,gray!40] (0,0) rectangle (14.2,8.3);
\begin{scope}[yshift=1cm]
\node at (7.1,6.7) {易错三：连续、两端等高，还必须检查内部可导};
\node[] at (3.8,6) {尖角：内部不可导};
\begin{scope}[shift={(3.8,3.9)},x=1.9cm,y=1.5cm]
\draw[->,gray!65] (-1.35,0)--(1.35,0) node[right] {$x$};
\draw[->,gray!65] (0,-1.45)--(0,0.8) node[above] {$y$};
\draw[blue!75!black,line width=1.15pt,,domain=-1:1,samples=160] plot (\x,{abs(\x)-1});
\fill[blue!75!black] (-1,0) circle (1.8pt);\node[above left] at (-1,0) {$-1$};
\fill[blue!75!black] (1,0) circle (1.8pt);\node[above right] at (1,0) {$1$};
\fill[orange!85!black] (0,-1) circle (1.8pt);
\end{scope}
\node[] at (3.8,1.3) {$|x|-1$};
\node[] at (10.3,6) {圆滑最低点：导数为零};
\begin{scope}[shift={(10.3,3.9)},x=1.9cm,y=1.5cm]
\draw[->,gray!65] (-1.35,0)--(1.35,0) node[right] {$x$};
\draw[->,gray!65] (0,-1.45)--(0,0.8) node[above] {$y$};
\draw[blue!75!black,line width=1.15pt,,domain=-1:1,samples=160] plot (\x,{(\x)*(\x)-1});
\fill[blue!75!black] (-1,0) circle (1.8pt);\node[above left] at (-1,0) {$-1$};
\fill[blue!75!black] (1,0) circle (1.8pt);\node[above right] at (1,0) {$1$};
\fill[orange!85!black] (0,-1) circle (1.8pt);
\draw[orange!85!black,line width=1.2pt] (-0.32,-1)--(0.32,-1);
\end{scope}
\node[] at (10.3,1.3) {$x^2-1$};
\end{scope}
\node[text width=13cm,align=center] at (7.1,.9) {同尺度对照：左图两侧斜率为 $-1,1$，在尖角处接不上同一个导数。};
\end{tikzpicture}
